% CX_GRH supplement, for the number-algorithm component.
% Finding ID CX-N13 agreed with /root/number_algorithms.
% Cite BIMS17 from koiran_bibliography.tex.

\subsection{CX-N13: nonresidues for primes dividing \texorpdfstring{$p-1$}{p-1} in supplied extensions}
\label{cx:n13}

\paragraph{Statement.}
Let $p$ be a prime, let $r$ be a specified prime divisor of $p-1$, and let
$F=\mathbb F_p[X]/(f)$ be supplied by an irreducible polynomial $f$
of degree $m$. An $r$-th nonresidue in $F$ can be constructed in
deterministic $\operatorname{poly}(m\log p)$ time under the stipulated
Dirichlet-GRH theorem $H_A$. In particular this holds for quadratic
nonresidues in every supplied extension of every odd prime field.
For $r=2$ and even $m$, the construction is already unconditional.

\paragraph{Established even-degree input.}
Bhargava--Ivanyos--Mittal--Saxena \cite[Corollary~1.3, PDF p.~2]{BIMS17}
gives deterministic $\operatorname{poly}(m\log p)$-time $r$-th-root
extraction when $r$ is constant and $r\mid\gcd(m,p-1)$.
The following remark and \S2 specify that the irreducible polynomial
defining the field is part of the input. For $r=2$, this covers every
even $m$ and every odd $p$. This is a prior unconditional theorem,
not a consequence newly obtained from either RH premise.
For variable $r\mid p-1$, their Theorem~1.2 also gives the transfer
once $\zeta_r$ is available and $r\mid m$: the supplied irreducible
$f$ satisfies their Property~1.1. We give the direct algebraic
argument for odd $r$, so the uniformity does not rely on a
constant-$r$ root-of-unity algorithm.

\begin{proof}[Proof of the all-degree transfer and output refinement]
Under $H_A$, CX-N02 supplies $a\in\mathbb F_p^\times$ with
$a^{(p-1)/r}\ne1$. Put $\zeta=a^{(p-1)/r}$, which has order $r$.
If $r\nmid m$, then, since $p\equiv1\pmod r$,
\[
 a^{(p^m-1)/r}
 =\zeta^{1+p+\cdots+p^{m-1}}
 =\zeta^m\ne1.
\]
Thus the constant $a$ is the required nonresidue in the supplied field.

Suppose next that $r\mid m$ and $r$ is odd. Set $Q=p^{m/r}$ and let
$\sigma:F\longrightarrow F$ be $u\longmapsto u^Q$.
The cyclic group $F^\times$ has order $(Q-1)S$, where
$S=1+Q+\cdots+Q^{r-1}$. The image of $u\mapsto u^{Q-1}$
has order $S$. Since $Q\equiv1\pmod r$, we have $r\mid S$,
so that image contains the element $\zeta$ of order $r$.
Consequently the $\zeta$-eigenspace of $\sigma$ is nonzero.
Compute a nonzero element $u$ of this eigenspace by solving
\[
 (\sigma-\zeta I)u=0
\]
as a linear system over $\mathbb F_p$ in the supplied polynomial
basis of $F$. The Frobenius matrix and a nonzero kernel vector can
be computed in polynomial bit complexity in $m\log p$.
Then $u^{Q-1}=\zeta$. Write $Q=1+rt$. Modulo $r^2$,
\[
 1+Q+\cdots+Q^{r-1}
 \equiv r+rt\frac{r(r-1)}2\equiv r\pmod{r^2},
\]
where oddness of $r$ is used in the second congruence. Therefore
\[
 u^{(p^m-1)/r}
 =\left(u^{Q-1}\right)^{(1+Q+\cdots+Q^{r-1})/r}
 =\zeta\ne1.
\]
This constructs the required $r$-th nonresidue. In this branch
$r\le m$; in the preceding branch CX-N02 and modular powering
are polynomial in $\log p$ independently of the numerical value
of $r$. The claimed uniform polynomial bound follows.

It remains to cover $r=2$ and even $m$, and here no RH assumption
is needed. Invoke the cited unconditional square-root algorithm.
Start with $z=-1$. While $z^{(p^m-1)/2}=1$, replace $z$ by any
square root supplied by that algorithm. Initially $z$ has order $2$.
A square root of an element of order $2^j$, for $j\ge1$, has order
$2^{j+1}$: squaring halves that order, and an element with square
of $2$-power order has no odd order component. Consequently each
iteration doubles the order. After at most
$v_2(p^m-1)-1\le\log_2(p^m)$ iterations, $z$ has the full
$2$-primary order of $F^\times$ and is a nonsquare. Euler tests and
the polynomially many square-root calls retain polynomial bit
complexity in $m\log p$. This branch does not call the initial
$H_A$-dependent prime-field search. The cases exhaust all $m$.
\end{proof}

This scope correction applies to the corresponding nonresidue
consequences of CX-N05 and CX-N06. It does not remove the
number-field Hecke hypothesis from Shoup's broader primitive-root
candidate-set theorem or supply the same reduction when
$r\mid p^m-1$ but $r\nmid p-1$. The extension model is supplied;
the statement does not claim an unconditional construction of that
model from $p$ and $m$ alone. In characteristic two every element
is a square, so a quadratic nonresidue does not exist.

An independent elementary verification of the even-degree quadratic
case, including the $p\equiv3\pmod4$ trace construction, is recorded
in the working proof notes. It is not presented here as a new
discovery of the established theorem.
